\documentclass[11pt]{article}
\usepackage[a4paper,margin=22mm]{geometry}
\usepackage{amsmath,amssymb,amsthm}
\usepackage{microtype}
\usepackage[colorlinks=true,linkcolor=blue,urlcolor=blue]{hyperref}
\newtheorem{theorem}{Theorem}
\newtheorem{corollary}[theorem]{Corollary}
\newtheorem{remark}[theorem]{Remark}
\title{Conditioning along the merging tree of a matrix product unitary}
\author{TNLean contributors}
\date{2 October 2026}
\begin{document}
\maketitle

\section{Question and source}

Styliaris, Trivedi, and Cirac construct quantum circuits for matrix product
unitaries by merging interval isometries along a binary tree
\cite{STC25}. Their Merging Lemma gives a depth estimate depending on the
conditioning of the cut used at each merge. Replacing every such cost by its
maximum gives their general depth estimate. This note retains the individual
costs and derives a sufficient condition for polynomial depth which allows
them to grow with the chain length.

The source is the local file
\texttt{references/2508.08160/main.tex}: the general conditioning parameter
and theorem appear at lines 1373--1402; the Merging Lemma at lines
2145--2156; the depth calculation at lines 2380--2422; and the outlook at
line 1405. The paper's uniform-bulk theorem already proves polynomial depth
under its full-rank assumption. Its general theorem applies to nonuniform
MPUs with an exponent depending on conditioning. Neither the estimates
below nor the source theorem settle polynomial circuit complexity for
arbitrary nonuniform MPUs of bounded bond dimension.

\section{Conditioning by level}

Fix a chain of length $N=2^L$, a leaf depth bound $a\geq0$, and an overhead
constant $C\geq0$. Let $q_j\geq1$ bound the cost of each merge at level $j$.
Suppose the circuit construction supplies depths satisfying
\begin{align}
 T_0&\leq a,&
 T_{j+1}&\leq q_j\bigl(2T_j+C2^{j+1}\bigr).
 \label{eq:recurrence}
\end{align}
For fixed physical and bond dimensions the local overhead in the source
can be bounded in this form. The constants must be chosen independently of
$N$ when an asymptotic statement about a family of MPUs is intended.
The quantities $q_j$ here are actual upper bounds usable in the recurrence;
an infimum over boundary choices must first be approximated by valid choices.

\begin{theorem}\label{thm:levels}
Put $Q_0=1$ and $Q_j=\prod_{i=0}^{j-1}q_i$. Then
\begin{align}
 T_L
 &\leq 2^LQ_L\left(a+C\sum_{j=0}^{L-1}Q_j^{-1}\right)
 \leq 2^LQ_L(a+CL).
 \label{eq:level-bound}
\end{align}
\end{theorem}

\begin{proof}
Divide (\ref{eq:recurrence}) by $2^{j+1}Q_{j+1}$ to obtain
\begin{align}
 \frac{T_{j+1}}{2^{j+1}Q_{j+1}}
 &\leq \frac{T_j}{2^jQ_j}+\frac{C}{Q_j}.
\end{align}
Summation gives the first inequality. Since $q_j\geq1$, we have
$Q_j\geq1$, which gives the second.
\end{proof}

\begin{corollary}\label{cor:polynomial}
If $Q_L\leq N^c$ for a constant $c\geq0$, then
\begin{align}
 T_L&\leq (a+C\log_2N)N^{c+1}.
\end{align}
Thus a polynomial product of conditioning costs suffices for polynomial
depth, without a constant bound on each cost.
\end{corollary}

For example, suppose all level costs are one except a single cost $N^c$.
Then $Q_L=N^c$, so Corollary~\ref{cor:polynomial} applies. Using only
$\max_jq_j=N^c$ instead would give an exponent growing like $\log N$.
This is an example of numerical cost data; an MPU realizing these costs is
not constructed here.

\section{Conditioning along branches}

The level maxima may still overestimate the depth when the most expensive
merges at different levels lie on different branches. A branch estimate
retains their positions.

Consider a perfect binary tree of height $L$. For a vertex $v$, let $h_v$
be its height and $n_v=2^{h_v}$ the number of leaves below it. Assign a
depth bound $D_v$ and, at each internal vertex, a conditioning cost
$\kappa_v\geq1$. Suppose each leaf has depth at most $a$ and the two
children $v_0,v_1$ of each internal vertex obey
\begin{align}
 D_v&\leq\kappa_v
       \bigl(2\max(D_{v_0},D_{v_1})+Cn_v\bigr).
 \label{eq:branch-recurrence}
\end{align}
The maximum reflects parallel implementation of the two child intervals.
This is an explicit hypothesis on the implementation: one must account
for its ancillary registers and geometric placement before applying it.

\begin{theorem}\label{thm:branches}
Let $W_v$ be the largest product of internal costs along a path from $v$
to a leaf, including the cost at $v$. An empty product is one. Then
\begin{align}
 D_v&\leq n_v(a+Ch_v)W_v.
 \label{eq:branch-bound}
\end{align}
If every branch product from the root is at most $N^c$, the root depth is
at most $(a+C\log_2N)N^{c+1}$.
\end{theorem}

\begin{proof}
The leaf case is the initial bound. For an internal vertex of height $h$,
write $M=\max(W_{v_0},W_{v_1})$. Then $M\geq1$,
$W_v=\kappa_vM$, and both children have $n_v/2$ leaves and height $h-1$.
The induction hypothesis and (\ref{eq:branch-recurrence}) give
\begin{align}
 D_v
 &\leq \kappa_v n_v\bigl((a+C(h-1))M+C\bigr)\\
 &\leq n_v(a+Ch)\kappa_vM.
\end{align}
This proves (\ref{eq:branch-bound}).
\end{proof}

\section{An obstruction to a universal conditioning bound}

There is a simple family for which rearranging the tree cannot give a
polynomial product of the paper's conditioning parameters.

\begin{theorem}\label{thm:parity}
Let $N\geq2$, let $Z=\operatorname{diag}(1,-1)$, and let $c,s$ be nonzero real
numbers with $c^2+s^2=1$. The unitary
\begin{align}
 U_N&=cI+i s Z^{\otimes N}
 \label{eq:parity-unitary}
\end{align}
has operator Schmidt rank two at every nontrivial cut, with normalized
Choi Schmidt coefficients $|c|,|s|$. At each cut, the conditioning parameter
defined in the source is
\begin{align}
 q_k&=\frac{1}{|cs|}.
 \label{eq:parity-conditioning}
\end{align}
If $c=\cos\theta$ and $s=\sin\theta$, this unitary has an exact circuit of
at most $2N-1$ nearest-neighbour gates, with no ancillary qubits.
\end{theorem}

\begin{proof}
The involution $P=Z^{\otimes N}$ is Hermitian, so
$(cI-isP)(cI+isP)=(c^2+s^2)I=I$. For a cut into two nonempty intervals,
write $P=P_A\otimes P_B$. Both $P_A$ and $P_B$ are traceless, and
\begin{align}
 U_N&=I_A\otimes cI_B+P_A\otimes isP_B.
\end{align}
On each interval the identity and parity operators are orthogonal in the
Hilbert--Schmidt inner product and have equal norms. After normalization
this is a Schmidt decomposition with coefficients $|c|,|s|$. It proves
rank two and supplies a minimal bond-two representation.

The source's sets of boundary Gram matrices are obtained by tracing the
left and right double layers against density operators $\sigma,\tau$
(\texttt{lem:isometries}, lines 1756--1818). For the displayed factorization,
put $x=\operatorname{tr}(\sigma P_A)$ and
$y=\operatorname{tr}(\tau P_B)$. These are real and range over $[-1,1]$.
Up to the harmless sign convention for the imaginary off-diagonal entries,
the Gram matrices are
\begin{align}
 G_A&=\begin{pmatrix}1&x\\x&1\end{pmatrix},&
 G_B&=\begin{pmatrix}c^2&icsy\\-icsy&s^2\end{pmatrix}.
 \label{eq:parity-grams}
\end{align}
They are positive definite precisely when $|x|<1$ and $|y|<1$.
Direct inversion of these two matrices gives
\begin{align}
 \operatorname{tr}\bigl[G_B^{-1}(G_A^{-1})^T\bigr]
 &=\frac{c^2+s^2}{c^2s^2(1-x^2)(1-y^2)}
 \geq\frac{1}{c^2s^2}.
 \label{eq:parity-gram-inverse}
\end{align}
Equality is attained at $x=y=0$, realized by maximally mixed density
operators. The infimum over these Gram sets, explicitly specified at line
2124 of the source, is therefore $q_k=1/|cs|$. The source
establishes invariance under invertible bond gauges at line 1379,
so the calculation does not require a prescribed canonical gauge.

Finally let $G_j$ be the controlled-NOT from qubit $j$ to qubit $j+1$ and
set $V=G_1G_2\cdots G_{N-1}$. Conjugation of the target $Z$ by a
controlled-NOT gives the product of the control and target $Z$ operators.
Applying this identity successively gives $VZ_NV^\dagger=Z^{\otimes N}$.
Consequently
\begin{align}
 U_N&=V e^{i\theta Z_N}V^\dagger.
\end{align}
There are $N-1$ gates in each sweep and one single-qubit rotation.
\end{proof}

For $\theta_N=2^{-N}$ we have
$|\sin\theta_N\cos\theta_N|\leq2^{-N}$, hence $q_k\geq2^N$ at every cut.
Any binary tree with $N=2^L$ leaves has a branch of length at least $L$,
so its largest branch product is at least $2^{NL}=N^N$.
Nevertheless the circuit above has linear depth. Thus bounded bond
dimension alone does not imply the polynomial branch-product hypothesis.
This example is not a counterexample to polynomial circuit complexity.
It shows that the conditioning of this particular construction need not
reflect the depth of a different implementation. Since $U_N$ approaches
the identity for these angles, it also gives no approximation lower bound.

The example does not contradict the uniform-bulk theorem: in that theorem
both the bulk and the boundary are fixed independently of $N$. Here the
angle, and hence the boundary coefficient, varies with $N$.

\begin{remark}[A larger class of boundary choices]
The Gram sets in the source are sufficient for interval isometries; they
need not contain every boundary choice with this property. For the parity
family, take the first tensor with components $I,Z$, each interior tensor
diagonal in the bond label with components $I,Z$, and the last tensor with
components $cI,isZ$. Choose $L_j=I$ at every internal left boundary and
$R_k=I/\sqrt2$ at every internal right boundary. Writing $P_{jk}$ for the
parity on $[j,k]$, the corresponding interval maps are
\begin{align}
 V_{1k}&=\frac{|0\rangle\otimes I+|1\rangle\otimes P_{1k}}{\sqrt2},
       &&k<N,\\
 V_{jk}&=\frac{|00\rangle\otimes I+|11\rangle\otimes P_{jk}}{\sqrt2},
       &&1<j\leq k<N,\\
 V_{jN}&=|0\rangle\otimes cI+|1\rangle\otimes isP_{jN},
       &&j>1.
\end{align}
Each is an isometry, by orthogonality of the displayed bond states and
$c^2+s^2=1$. The full interval is $U_N$. At every merge these choices give
$\sqrt{\operatorname{tr}[R_k^{-2}(L_{k+1}^{-2})^T]}=2$.
When $c^2\ne s^2$, $R_k^2=I/2$ lies outside the source's right Gram set:
its diagonal is not $(c^2,s^2)$. Thus the exponential value of $q_k$ does
not obstruct a constant merge cost after enlarging the permitted boundary
class. The displayed choices are simultaneous for all intervals, which is
necessary for recursively applying the merging construction.
\end{remark}

For this family $s_{\min}=\min(|c|,|s|)$, and
\begin{align}
 \frac{1}{s_{\min}}&\leq q_k\leq\frac{\sqrt2}{s_{\min}}.
\end{align}
Thus the source's estimate using the smallest Schmidt coefficient has the
correct dependence for its conditioning parameter, even in a family with
an elementary linear-depth implementation.

\section{Product-projection phase gates}

The same enlargement works for a family with projection-valued local
factors. Let $N\geq2$ and let $P_1,\ldots,P_N$ be orthogonal projections, not necessarily
identical, and let $z\in\mathbb C$ satisfy $|z|=1$. Put
$P_{jk}=\bigotimes_{t=j}^k P_t$. The phase gate
\begin{align}
 U_N&=I+(z-1)P_{1N}
\end{align}
is unitary, since it acts by $z$ on the range of $P_{1N}$ and by one on its
orthogonal complement. It has an open-boundary representation of bond at
most two: the first tensor has components $I,P_1$, the interior tensors are
diagonal in the bond label with components $I,P_j$, and the last tensor
has components $I,(z-1)P_N$. It is a specialization of the MPO-projector
unitaries described at line 817 of the source.

\begin{theorem}\label{thm:projection-frames}
The product-projection phase gate admits simultaneous interval isometries
with bond contraction norm two, independently of $z$, $N$, and the local
projections.
\end{theorem}

\begin{proof}
Choose the left and right bond vectors, respectively, as the columns of
\begin{align}
 C_L&=\begin{pmatrix}1&1/2\\0&1/2\end{pmatrix},&
 C_R&=\begin{pmatrix}1&-1\\0&1\end{pmatrix}.
\end{align}
For a projection $P$ on an interval, the prefix, interior, and suffix maps
have components
\begin{align}
 V_{\mathrm{pre}}(P)&=\begin{pmatrix}I-P\\P\end{pmatrix},\\
 V_{\mathrm{int}}(P)&=\begin{pmatrix}I-P/2\\P/2\\-P/2\\P/2\end{pmatrix},\\
 V_{\mathrm{suf}}(P,z)&=\begin{pmatrix}I+(z-1)P/2\\(z-1)P/2\end{pmatrix}.
 \label{eq:projection-isometries}
\end{align}
The order of the four interior components is the lexicographic order of
the left and right bond labels. Using $P^\dagger=P$ and $P^2=P$, the prefix
and interior Gram products are $I$. For the suffix, the coefficient of $P$
in its Gram product is
\begin{align}
 \frac{z-1}{2}+\frac{\bar z-1}{2}
 +\frac{|z-1|^2}{2}&=\frac{|z|^2-1}{2}=0.
\end{align}
Its Gram product is therefore also $I$.

Let the contraction of a right and left bond state have coefficient matrix
\begin{align}
 K&=\begin{pmatrix}1&-1\\1&1\end{pmatrix}.
\end{align}
Direct multiplication gives
\begin{align}
 C_R^T K C_L&=I,& \operatorname{tr}(K^\dagger K)&=4.
 \label{eq:projection-contraction}
\end{align}
The first identity contracts the bond columns with the Kronecker delta.
Thus joining adjacent interval maps multiplies matching physical factors,
$P_{jk}\otimes P_{k+1,l}=P_{jl}$, and gives exactly the parent interval map,
including the prefix, suffix, and full-chain cases. The second identity is
the squared norm of the contraction functional. Both frames are invertible,
with determinants $1/2$ and one, and their Gram matrices satisfy
\begin{align}
 \operatorname{tr}\bigl[(C_R^\dagger C_R)^{-1}
                ((C_L^\dagger C_L)^{-1})^T\bigr]&=4.
\end{align}
The isometry and contraction identities hold simultaneously for all
intervals, rather than requiring a new boundary choice at each merge.
\end{proof}

The amplification argument of the source's Merging Lemma uses precisely
interval isometries and a norm bound for their contraction. The provenance
of the boundary matrices in exterior density operators is used to obtain
isometries, and is not needed once these identities are supplied directly.
The frames in Theorem~\ref{thm:projection-frames} therefore give a
phase-independent merging recurrence with conditioning cost two. For a
balanced tree and fixed physical dimensions this yields the polynomial
bound from Theorem~\ref{thm:levels}. This is an improvement of the
conditioning certificate for this family, not a claim of a new optimal
circuit for product-controlled phase gates.

For comparison, suppose each local projection is nonzero and proper and
$z\ne1$. At a nontrivial cut, the exterior density-operator Gram matrices
in the source are
\begin{align}
 G_A&=\begin{pmatrix}1&x\\x&x\end{pmatrix},&
 G_B&=\begin{pmatrix}1&(z-1)y\\(\bar z-1)y&|z-1|^2y\end{pmatrix},
\end{align}
where $0<x,y<1$. Every such expectation is attained by a suitable density
operator. The identity $|z-1|^2+(z-1)+(\bar z-1)=0$ gives
\begin{align}
 \operatorname{tr}\bigl[G_B^{-1}(G_A^{-1})^T\bigr]
 &=\frac{1}{|z-1|^2x(1-x)y(1-y)}.
\end{align}
Minimization at $x=y=1/2$ gives the source parameter
$q_k=4/|z-1|$. Its divergence as $z\to1$ is absent from the enlarged-boundary
certificate, whose cost remains two.

\section{Choosing the interval tree}

For an ordered chain, the same Merging Lemma permits a quantitative search
over trees. Suppose $\kappa_k$ is a valid cost at cut $k$, and $b_{i,j}$
is a specified nonnegative overhead for merging the interval $[i,j]$.
Define recursively
\begin{align}
 F(i,i)&=a_i,\\
 F(i,j)&=\min_{i\leq k<j}\kappa_k
     \bigl(2\max(F(i,k),F(k+1,j))+b_{i,j}\bigr).
 \label{eq:interval-minimum}
\end{align}
There are finitely many cuts, so each minimum is attained. Induction on
interval length shows that $F(i,j)$ is the least upper bound obtained by
this recurrence among all ordered binary merging trees on $[i,j]$.
It is an attainable bound within the assumed merging construction, not the
minimum circuit depth among arbitrary implementations. Evaluating all
intervals and their cuts uses $O(N^3)$ arithmetic operations when the costs
are already supplied. No claim about computing these costs from a tensor
is needed for this statement.

Theorem~\ref{thm:parity} shows that a polynomial branch-product bound cannot
hold for every bounded-bond family. Tree optimization can still improve
the source construction when the difficult cuts are distributed unevenly.
A proof of polynomial complexity for all MPUs must also handle families
for which every tree has large conditioning products.

\section{Two points in the source calculation}

\begin{remark}[Maximum over cuts]
The main text defines $q=\max_kq_k$ at line 1382. The appendix writes
$q=\min_kq_k$ at line 2123. A common bound for every merge requires the
maximum, as in the main text. Retaining the individual costs avoids this
ambiguity. This observation concerns the displayed definitions, not a
counterexample to the main theorem.
\end{remark}

\begin{remark}[The endpoint $q=1$]
For constant $q>1$, the first estimate in (\ref{eq:level-bound}) yields
\begin{align}
 T_L&\leq (2q)^L
      \left(a+C\frac{q}{q-1}(1-q^{-L})\right).
\end{align}
For fixed $q>1$ this recovers the source's exponent. At $q=1$, the
recurrence permits $T_L=2^L(a+CL)$; with $a=0$, the sequence $CL2^L$
satisfies it with equality. This numerical observation gives no circuit
lower bound.

For the source's density-operator Gram sets, the endpoint has a further
restriction. Put $S=R_k^2$ and $H=(L_{k+1}^2)^T$, of bond dimension $r$.
Unitarity gives $\operatorname{tr}(SH)=1$. The positive eigenvalues
$\lambda_i$ of $S^{1/2}HS^{1/2}$ consequently sum to one, and
\begin{align}
 \operatorname{tr}(S^{-1}H^{-1})
 &=\sum_{i=1}^r\lambda_i^{-1}\geq r^2.
\end{align}
Cauchy--Schwarz gives $q_k\geq r$. Thus unit conditioning at every cut
forces a minimal bond-one representation, hence an on-site product unitary.
The numerical endpoint does not contradict the source's circuit theorem.
\end{remark}

\section{Formalization and remaining questions}

The module \texttt{TNLean/MPS/MPU/TreeDepthBounds.lean} proves the second
inequality of (\ref{eq:level-bound}), its polynomial corollary for a natural
exponent, the constant-cost specialization, and the endpoint identity.
The module \texttt{TNLean/MPS/MPU/ParityConditioning.lean} proves the
algebraic identity in (\ref{eq:parity-gram-inverse}). The module
\texttt{TNLean/MPS/MPU/ProjectionPhaseIsometries.lean} proves the three
isometry identities in (\ref{eq:projection-isometries}), the identities in
(\ref{eq:projection-contraction}), invertibility of the frames, and their
inverse-Gram trace. The sharper summation and branch estimates, density-state
minimizations, complete tensor representations and circuits, and interval
optimization remain mathematical proofs rather than Lean theorems.

The source leaves measurement-assisted depth improvements and
higher-dimensional extensions open. Polynomial depth for all nonuniform
MPUs of bounded bond dimension also remains unresolved. The examples show
that an affirmative proof must avoid a universal polynomial bound on the
source's conditioning products. Enlarging the boundary class is one possible
approach, but controlled compatible frames for general tensors remain to be
established.

\begin{thebibliography}{1}
\bibitem{STC25}
G.~Styliaris, R.~Trivedi, and J.~I.~Cirac,
\emph{Quantum Circuits for Matrix-Product Unitaries},
Physical Review Letters \textbf{135}, 260602 (2025),
\href{https://arxiv.org/abs/2508.08160}{arXiv:2508.08160v2}.
\end{thebibliography}
\end{document}
